% Generated draft: requires mathematical and source review.
\documentclass[11pt]{article}
\usepackage[margin=25mm]{geometry}
\usepackage{fontspec}
\setmainfont{Latin Modern Roman}
\usepackage{amsmath,amssymb,mathtools,unicode-math}
\setmathfont{Latin Modern Math}
\usepackage{enumitem,needspace,xurl,hyperref}
\setlength{\parindent}{0pt}
\setlength{\parskip}{5pt}
\setlength{\emergencystretch}{3em}
\newcommand{\N}{\mathbb{N}}
\newcommand{\Z}{\mathbb{Z}}
\newcommand{\Q}{\mathbb{Q}}
\newcommand{\R}{\mathbb{R}}
\newcommand{\C}{\mathbb{C}}
\newcommand{\F}{\mathbb{F}}
\newcommand{\A}{\mathbb{A}}
\newcommand{\PP}{\mathbb{P}}
\newcommand{\E}{\mathbb{E}}
\newcommand{\eps}{\varepsilon}
\newcommand{\dd}{\,\mathrm{d}}
\begin{document}
\section*{UnsolvedMath ID: 30004690}
\textbf{Source problem code:} catalogue entry
\subsection{Definitions and mathematical statement}
Let \(\mathbb{P}^{1}\) denote the complex projective line, let \(D=\{z\in\C:|z|<1\}\), and set \(X=\mathbb{P}^{1}\times D\).

The requested construction is a solution \(u\) of the homogeneous complex Monge--Amp\`ere equation on \(X\) for which there exists an interior point \(p\in X\) such that \(u\) is not of class \(C^{2}\) on any open neighborhood of \(p\). Here \(C^{2}\) means twice continuously differentiable in local real coordinates. Thus the requested failure of regularity must occur in the interior, rather than solely at the boundary \(\mathbb{P}^{1}\times\partial D\).

\textbf{Unresolved formulation qualifications.} The supplied record does not specify the precise Monge--Amp\`ere operator, any background differential form, the admissible class of solutions, or the weak interpretation of the equation needed for a solution that is not \(C^{2}\). It also supplies no boundary data or regularity requirements for such data. Consequently, this is the stated construction target, but a fully specified equation and boundary-value problem cannot be recovered from the supplied evidence. In particular, the record does not determine whether smooth boundary data are required.
\subsection{Short English statement}
Construct a solution of the intended homogeneous complex Monge--Amp\`ere problem on the complex projective line times the unit disc that fails to be twice continuously differentiable near at least one interior point. The equation, solution class, and boundary requirements need clarification before this becomes a fully precise construction problem.
\subsection{Sources}
\begin{itemize}
\item Local UnsolvedMath registry, supplied record with permanent ID \(30004690\), title \emph{Nonsmooth Homogeneous Complex Monge--Amp\`ere Solutions}, exact source problem code \texttt{catalogue entry}. The record belongs to \emph{Oberwolfach Reports Open Problems}, gives proposed year \(2021\), and names no individual proposer.
\item Supplied source citation: \emph{Geometric Methods of Complex Analysis (2022)}, Oberwolfach Reports, pp.~1291--1345, DOI \texttt{10.4171/owr/2021/25}. Supplied source and literature link: \url{https://doi.org/10.4171/owr/2021/25}. The background describes these reports as records of mathematical workshops and their open-problem discussions.
\item Supplied literature reference labeled ``Ross--Witt Nystr\"om, maximal rank'': \url{https://doi.org/10.2140/apde.2019.12.493}. No fuller bibliographic citation is supplied.
\item Literature assessment dated \texttt{2026-08-22}: the supplied curation labels the problem open and states that the 2021 report leaves unresolved the explicit construction of a solution outside \(C^{2}\) on \(\mathbb{P}^{1}\times D\); it further states that later work constructs solutions of non-maximal rank but does not determine their precise failure of regularity. This is supplied dated curation, not an independently verified literature assessment. The supplied background explicitly requests independent verification of time-sensitive status and cited claims.
\end{itemize}
\end{document}
